\section{Introduction}

Human-in-the-loop reinforcement learning (HIL-RL) is attractive for robotic manipulation because human interventions can redirect an agent away from unsafe, unproductive, or low-value exploration while adding corrective data to the learning process. Recent robot-learning systems combine interventions, demonstrations, and off-policy updates to improve sample efficiency in manipulation settings \cite{luo2025hilserl,ball2023rlpd}. More broadly, surveys of human involvement in learning systems emphasize that the human can shape development, training, evaluation, and deployment rather than merely supply an initial dataset \cite{retzlaff2024hitl}. In this setting, intervention timing is not only a control-design problem. It is a human-attention allocation problem: the trigger decides when a scarce, costly, and potentially safety-critical human resource is spent.

This allocation view changes what must be measured. A learned trigger can look human-efficient if it is compared with a single random baseline whose budget is too high or too low. A checkpoint can look strong if it is selected from a noisy training curve and evaluated only once. A trigger score can appear meaningful at the episode level while still spending budget at the wrong training phase or with poor calibration. These failure modes matter for robotics because real corrections carry opportunity cost, operator delay, interface burden, and safety implications. An intervention policy that merely moves cost into poorly measured parts of training is not yet evidence of better human-attention allocation.

FORESIGHT-HIL began from a method-positive hypothesis: a learning-value value-of-information trigger might request human help more efficiently than random intervention. The current evidence does not support that hypothesis. In the Lift case study, the LV-VoI scale3 trigger initially appeared plausible when compared with a higher-cost random baseline, but a random-budget sweep changed the conclusion. The cost-matched random b350 baseline reached 439/500 repeated successes (87.8\%) at 177.0 best-step human steps, while LV-VoI scale3 reached 416/500 repeated successes (83.2\%) at 202.0 human steps. Under the current evidence registry, the learned trigger is therefore dominated rather than Pareto-superior.

The paper turns this reversal into a diagnostic protocol for human-attention allocation claims in HIL-RL. The protocol asks whether a trigger survives three validity gates: random baselines swept around the trigger's realized human cost, repeated checkpoint evaluation with binomial uncertainty, and trace-level diagnostics of when and where interventions start. These checks are deliberately modest. They do not require a new robot platform or a full benchmark suite, but they do require enough measurement discipline to prevent an allocation claim from resting on a mismatched budget, a lucky checkpoint, or an uninterpreted aggregate success rate.

The resulting contribution is bounded but useful for automation science. We present a robotic manipulation case study in robosuite \cite{robosuite2020} using an SAC-style off-policy learner \cite{haarnoja2018sac,sac2018applications} and a scripted privileged-state oracle as the simulated intervention source. We show that cost matching reverses the apparent LV-VoI advantage, that two lightweight repairs remain dominated by same-seed random baselines, and that intervention traces point to budget-control and score-timing mismatch rather than a simple far-from-object failure. The paper therefore does not claim that the current LV-VoI trigger is superior. Its claim is that HIL-RL trigger papers should pass cost-matched, repeated-evaluation, and trace-level diagnostic gates before presenting intervention timing as a human-efficient improvement.
